% 第5章图：Bartels--Stewart 逐列求解 TY+YS=D 的顺序与依赖（上三角情形，d=m=5，当前列 j=3）。
% 第 j 列方程：(T+s_jj I) y_j = d_j - sum_{l<j} s_{l,j} y_l；列从左到右，列内从下到上回代。
\begin{tikzpicture}[x=1mm,y=1mm,>=Stealth,
  cellbox/.style={draw=FigureStroke,line width=.5pt},
  known/.style={cellbox,fill=FigureGreenFill},
  current/.style={cellbox,fill=FigureYellowFill},
  head/.style={font=\figurefont\fontsize{8}{10}\selectfont},
  note/.style={font=\FigureNoteFont\fontsize{7.5}{9}\selectfont,text=FigureMuted}]
\def\c{4}
% ---------- 面板一：S 的上三角格局，第 j=3 列 ----------
\foreach \i in {1,...,5}{\foreach \j in {\i,...,5}{
  \pgfmathsetmacro{\xa}{(\j-1)*\c}\pgfmathsetmacro{\ya}{-(\i-1)*\c}
  \ifnum\j=3
    \ifnum\i<3 \draw[known] (\xa,\ya) rectangle ++(\c,-\c);
    \else \draw[current] (\xa,\ya) rectangle ++(\c,-\c); \fi
  \else \draw[cellbox,draw=FigureGrid] (\xa,\ya) rectangle ++(\c,-\c); \fi
}}
\draw[rounded corners=1.8pt,FigureStroke,line width=.7pt] (0,0) rectangle (20,-20);
\node[head] at (10,6.5) {$\bm S$};
\node[note] at (10,2.2) {第$j$列};
\node[note,align=center] at (10,-25.5) {$s_{\ell,j}\,(\ell<j)$带入已求列\\$s_{j,j}$并入当前系数};
% ---------- 面板二：Y 的列，左到右 ----------
\foreach \j/\st in {1/known,2/known,3/current,4/cellbox,5/cellbox}{
  \draw[\st] ({32+(\j-1)*\c},0) rectangle ++(\c,-20);
  \node[font=\figurefont\fontsize{7.5}{9}\selectfont] at ({32+(\j-0.5)*\c},-10) {$\bm y_{\j}$};
}
\draw[rounded corners=1.8pt,FigureStroke,line width=.7pt] (32,0) rectangle (52,-20);
\node[head] at (42,6.5) {$\bm Y$};
\draw[->,FigureStroke,line width=.7pt] (32,2.2)--(52,2.2);
\node[note,align=center] at (42,-25.5) {绿色已求，黄色当前\\按$j=1,\ldots,d$从左到右};
% ---------- 面板三：当前列的上三角回代 ----------
\foreach \i in {1,...,5}{\foreach \j in {\i,...,5}{
  \pgfmathsetmacro{\xa}{64+(\j-1)*\c}\pgfmathsetmacro{\ya}{-(\i-1)*\c}
  \ifnum\i=\j \draw[current] (\xa,\ya) rectangle ++(\c,-\c);
  \else \draw[cellbox,draw=FigureGrid] (\xa,\ya) rectangle ++(\c,-\c); \fi
}}
\draw[rounded corners=1.8pt,FigureStroke,line width=.7pt] (64,0) rectangle (84,-20);
\node[head] at (74,6.5) {$\bm T+s_{j,j}\bm I$};
\draw[current] (87,0) rectangle (91,-20);
\node[font=\figurefont\fontsize{7.5}{9}\selectfont] at (89,-10) {$\bm y_j$};
\draw[->,FigureGreen,line width=1pt] (93.5,-19)--(93.5,-1);
\node[head] at (97,-10) {$=$};
\draw[cellbox,fill=white] (100,0) rectangle (104,-20);
\node[font=\figurefont\fontsize{7.5}{9}\selectfont] at (102,-10) {$\bm r_j$};
\node[note,align=center] at (84,-25.5) {$\bm r_j=\bm d_j-\sum_{\ell<j}s_{\ell,j}\bm y_\ell$\\对角元$t_{i,i}+s_{j,j}\neq0$，从下到上回代};
\end{tikzpicture}
